\section{Shared response libraries}\label{sec:21.6}

At $|\vect{q}| = \alpha m_e$ and $\omega = \SI{10}{eV}$, the screened structure factor of silicon with local fields is \SI{4.2e7}{eV^2} in the published random-phase tables of Knapen, Kozaczuk and Lin~\cite{knapen2021dark} and \SI{4.9e7}{eV^2} in the empirical-pseudopotential model of \cref{we:8-pairs}. Both calculations are at the same point of the truncation label, $\mathcal{X} = \mathrm{RPA + LF}$. They agree on the independent-particle response to 3 percent and differ in the screening, which leaves 0.24 of the unscreened value in the tables and 0.29 in the model (\cref{we:21-libentry}). At \SI{3}{MeV} two tabulations of the same independent-particle response give rates that differ by a third (\cref{sec:10.10}). A limit can therefore be reproduced only from the table it was computed with, and recast for another coupling, halo or readout only if that table holds the rows, directions and final states the new calculation needs. The schema of \cref{tab:21-schema} lists what a response library must tabulate for a recast, which conventions each table must state, and which section of \crefrange{chap:23}{chap:32} specifies the calculation that fills it.

\subsection{Tables of a response library}\label{sec:21.6-tables}

The axes of a library follow from the rate kernel \cref{eq:21.1-kernel}. The velocity of the dark matter enters it in two places: in the coupling matrix $\mathfrak{D}_{ab}(\vect{q},\vect{v})$ and in the energy $\omega_{\vect{q}}$ fixed by energy conservation. The response $S_{ab}(\vect{q},\omega;\Lambda)$ depends on $\vect{q}$, $\omega$ and the final-state class, and $P_{\mathrm{casc}}(N_{\mathrm{eh}}|\Lambda)$ on the class alone. A library therefore has no velocity axis and holds nothing about the halo. It tabulates functions of $|\vect{q}|$, $\hat{\vect{q}}$ and $\omega$, labeled by a pair of coupling rows $(a,b)$, a class $\Lambda$ and a pair number. The coupling model enters through $\mathfrak{D}$, built from the operator coefficients of \cref{tab:8-operators} and contracted by the user as $\sum_{ab}\mathfrak{D}_{ab}S_{ab}$, so the response is stored uncontracted, one function per row pair. Velocity-dependent couplings need no velocity axis either: their kernels are polynomials in the transverse velocity, and after the orientation average every power is fixed by $\eta(\vmin)$ (\cref{eq:21.1-eta2}).

The electronic response is the largest part. For silicon and germanium the minimal table of \cref{tab:10-minimal} reduces every nonrelativistic operator to the 12 real functions per $(\vect{q},\omega)$ of \cref{app:E.16}, with four more in each version for operator $\mathcal{O}_2$. A crystal without the combined symmetry of inversion and time reversal, such as GaAs, keeps the imaginary parts of the cross terms as well. On the grid of \cref{app:E.16}, 250 values of $|\vect{q}|$ up to $25\,\alpha m_e$, 30 directions and 1000 energies up to \SI{50}{eV}, the 12 functions fill $7.5\times10^6$ nodes, \SI{720}{MB} in double precision. A screened response at a higher level, a time-dependent density-functional kernel or the Bethe--Salpeter equation, needs the full independent-particle matrix $R^0_{\vect{G}\vect{G}'}(\vect{q},\omega)$, which is too large to store on the same grid and is kept on a coarse subset of the nodes.

The direction of $\vect{q}$ is kept as far as the analyses that use it require. A speed-only kernel, \cref{eq:21.1-gammav}, needs the direction average $\bar S_{ab}$. The daily modulation (\cref{sec:11.7}) and the directional halo-independent bound of \cref{eq:21.2-multipole} need the dependence on $\hat{\vect{q}}$ in crystal axes. In a cubic crystal the direction dependence is carried, away from the edges, by the cubic-harmonic components $S_\ell/S_0$ for $\ell = 4$, 6 and 8. Near the zero-phonon edge the response changes on a scale of \SI{0.1}{rad} in direction (\cref{sec:13.9}), and there the library keeps the full grid of directions. The lattice tables of \cref{app:E.17} hold the response per cell, the probability $P_{\mathrm{ph}}(n|\vect{q},\omega)$ that $n$ phonons were created and the coherence correction, on a grid of $1.2\times10^7$ values per quantity.

The energy grid has to resolve the edges. The power law of the rate near an edge, \cref{eq:21.3-power}, rests on the onset of the response. The computed phonon-assisted loss function of silicon follows its $x^3$ onset only for the first \SI{1.3}{meV} above the onset, and the zero-phonon continuum follows its onset for \SI{49}{meV} (\cref{sec:21.3-power}). A step of \SI{0.05}{eV}, adequate for the bulk of the spectrum, has no node inside the phonon-assisted range, which is 38 times narrower. The library therefore adds a sub-grid with steps of a fraction of a meV within some \SI{50}{meV} of each onset, and lists for each channel the derived edge entries of \cref{tab:21-edge}: the edge mass $m_{\mathrm{edge}}$, the onset energy and the onset exponent $p_{\mathrm{on}}$.

The final-state tables carry the per-pair weights of \cref{eq:8.7-screened}. A full list is out of reach: the library grid has $10^3$ points $\vect{k}$ and $4\times56$ band pairs in silicon, $2.24\times10^5$ pairs per node and $1.7\times10^{12}$ numbers, \SI{13}{TB} for each coupling row. The cascade does not need the list. In the energy-only reduction of \cref{sec:18.3} it needs the seed distribution $P_{\mathrm{seed}}(E_{\mathrm{e}},E_{\mathrm{h}}|\vect{q},\omega,\vect{v})$ of \cref{eq:18.3-seed}, which for one coupling row depends on neither the velocity nor the model. In the zero-phonon sector $E_{\mathrm{e}} + E_{\mathrm{h}} = \omega - \Eg$, so one variable, the electron share $f_e$, describes the seed, and 20 bins of it take \SI{1.2}{GB} per row pair on the full grid. A seed with a phonon adds the phonon energy, since then $E_{\mathrm{e}} + E_{\mathrm{h}} = \omega - \omega_{\vect{Q}} - \Eg$, and a valley label for the electron keeps the momentum memory of \cref{sec:18.1}. The weights are the retarded ones of \cref{sec:8.7}, which label the late-time state the cascade receives; in silicon and germanium they equal the advanced weights pair by pair, and in GaAs the retarded and advanced weights differ by 15 percent of the total. The split between the zero-phonon and the phonon-dressed classes follows either the off-shell window or the Dyson convention of \cref{sec:8.6}, and the totals are independent of the choice only when one convention is declared for every table of the library.

The cascade tables hold $P_{\mathrm{casc}}(N_{\mathrm{eh}}|\Lambda)$ for each cascade model and parameter set, through the subtree generating functions $\mathfrak{G}_{\mathrm{e}}(x;s)$ and $\mathfrak{G}_{\mathrm{h}}(x;s)$ of \cref{eq:18.3-pcasc} on a grid of carrier energies $x$. \Cref{app:E.24} specifies that grid, a step of about $\Eg/50$ up to $20\,\Eg$ and the asymptotic Gaussian beyond, with its acceptance tests. The same tables carry the marks $\mathfrak{m}$ of \cref{eq:18.3-obs} that a readout needs: the phonon spectrum, the separate electron and hole counts, the carrier positions and the stored energy. The detector response $P_{\mathrm{det}}(\mathsf{R}|N_{\mathrm{eh}},\mathfrak{m})$ of \cref{eq:18.3-obs} belongs to the experiment, which measures and publishes it with its collection efficiencies (\cref{chap:20}), and the velocity distribution belongs to the analysis, which contracts it with the kernel in \cref{eq:21.1-functional}.

The uncertainty tables make the theory band of a limit reproducible. For each truncation $k$ active at a node, the library stores the next-order table $S^{\mathfrak{T}_k^+}$ itself rather than the ratio $\delta_k$ of \cref{eq:8.8-delta}. The matrix form \cref{eq:8.8-geodesic} is built from $S^{\mathfrak{T}_k^+}$, and a ratio is undefined wherever $S^{\mathfrak{T}}$ vanishes, which is where a new-support template acts (\cref{eq:8.8-open}). Beside the next-order table stand the type of the template in the classes of \cref{sec:8.8} (normalization, tilt, threshold or new support), its basis (computed, model, estimate) and a key naming its parameter $\vartheta_k$ in the likelihood of \cref{sec:21.4}. The key enforces coherence: one $\vartheta_k$ acts on every node, class, target and experiment that shares the approximation (\cref{sec:8.8}), and two tables carrying the same key are varied together. A template known only as an inequality is flagged as a bound; one example is the tabulation spread of \cref{sec:10.10}, which bounds the width of a spread that enters a likelihood with that width (\cref{sec:21.4-like}). A node with no template carries a flag that names the missing calculation. Energies at which no central value can be defended are flagged as excluded; at the superconducting edge of an aluminum film (\cref{sec:15.2}) these are the energies below $3\Delta_{\mathrm{BCS}}$. The band of \cref{sec:21.4} is then built with these omissions named rather than with invented widths. A second calculation at the same truncation label is stored as an alternative table; the spread between the two is the uncertainty of the approximations that the label leaves out, as for the two silicon tabulations of \cref{sec:10.10}.

The derived tables are those that analyses would otherwise recompute for every search. The step column $\mathcal{A}_i(v)$ of \cref{eq:21.2-column}, per bin and cascade model, turns the halo-independent bound of \cref{eq:21.2-hib} into a minimum over a few numbers, and it follows from the response by cumulative sums over $\vmin$. It is also a check: for a nonnegative kernel it never decreases with $v$ (\cref{eq:D.46-monotone}), and a decrease flags either a sign error in a template or the final-state leakage of \cref{eq:21.2-fsleak}, at most $0.035\,\varrho^{\mathrm{leak}}_i$ per silicon pair bin. The direction-resolved stream kernels of \cref{app:E.27} complete the set for directional searches.

\begin{table}[tbp]
  \caption{Schema of a response library for one material. Every table carries the truncation label $\mathfrak{T}$ of \cref{eq:8.5-label} with the metadata and conventions of \cref{sec:21.6-conventions}, a content hash and the list of its parents (\cref{sec:21.6-versions}). Energies $\omega$ are given to the target, momenta $\vect{q}$ point from the dark matter to the target and are not reduced to the first zone. The last column names the sections that specify the calculation filling the table.}
  \label{tab:21-schema}
  \centering\footnotesize
  \begin{tabular}{@{}>{\raggedright\arraybackslash}p{2.1cm}>{\raggedright\arraybackslash}p{2.8cm}>{\raggedright\arraybackslash}p{6.6cm}>{\raggedright\arraybackslash}p{2.5cm}@{}}
    \toprule
    Table & Axes and labels & Content & Specified in \\
    \midrule
    Material inputs & --- & structure and measured lattice constant; the gap of the scissor and the gap used by the cascade, separately; functional, pseudopotentials, cutoffs, grids, number of bands; force constants of the shared harmonic base & \cref{app:E.1,app:E.17} \\
    Electronic response & $|\vect{q}|$, $\hat{\vect{q}}$, $\omega$; rows $(a,b)$ & minimal table of \cref{tab:10-minimal} (12 real functions in Si and Ge, 20 with $\mathcal{O}_2$), screened and independent-particle, per cell; $R^0_{\vect{G}\vect{G}'}$ on a coarse subset of nodes & \cref{app:E.6,app:E.7,app:E.16} \\
    Directional components & $|\vect{q}|$, $\omega$; $\ell$ & $S_\ell/S_0$ for $\ell = 4$, 6, 8; full direction grid within the edge sub-grid & \cref{app:E.16} \\
    Lattice response & $|\vect{q}|$, $\hat{\vect{q}}$, $\omega$; $T$ & response per cell; $P_{\mathrm{ph}}(n|\vect{q},\omega)$; branch content; coherence correction & \cref{app:E.17} \\
    Further channels & as the parent table & phonon-assisted loss function; exciton lines and core shells; joint response of the Migdal effect; spin-wave, superconducting and Dirac responses & \cref{app:E.14,app:E.15,app:E.18,app:E.20} \\
    Final-state classes & $|\vect{q}|$, $\omega$, rows; $f_e$, $\omega_{\vect{Q}}$, valley & seed distribution $P_{\mathrm{seed}}$ per class from retarded pair weights; convention of the zero-phonon/phonon-dressed split & \cref{app:E.6,app:E.7} \\
    Cascade & carrier energy $x$; model and parameters & $\mathfrak{G}_{\mathrm{e}}$, $\mathfrak{G}_{\mathrm{h}}$ and $P_{\mathrm{casc}}(N_{\mathrm{eh}}|\Lambda)$; marks $\mathfrak{m}$ for the readout & \cref{app:E.13,app:E.24} \\
    Theory uncertainties & as the parent table; key of $\vartheta_k$ & next-order table $S^{\mathfrak{T}_k^+}$; type and basis of the template; flags for bounds, missing templates and excluded energies; alternative tables at the same $\mathfrak{T}$ & as the parent table \\
    Checks & as the parent table & $f$-sum fraction; smallest eigenvalue of $S_{ab}$; Kramers--Kronig residual; detailed balance; symmetry-forbidden parts; cascade normalization, moments and the single-pair limit & as the parent table \\
    Derived & bin $i$, speed $v$; channel & step columns $\mathcal{A}_i(v)$; edge mass, onset energy, $p_{\mathrm{on}}$; direction-resolved stream kernels & \cref{app:E.27} \\
    \bottomrule
  \end{tabular}
\end{table}

\subsection{Conventions declared with each table}\label{sec:21.6-conventions}

A number in a table is usable only with the conventions that turn it into a rate, and each convention that a user has to guess is a possible error of a factor of order one. The momentum $\vect{q}$ points from the dark matter to the target and $\omega$ is the energy given to the target, negative at finite temperature for a de-excitation; detailed balance, $S(\vect{q},-\omega) = \ee^{-\beta\omega}S(\vect{q},\omega)$, relates the two signs, and the table states the temperature it holds. The temperature is metadata, distinct from the thermal order $n_T$ of the label, which counts powers of the phonon occupations (\cref{sec:8.5}).

Next comes the normalization, the most common source of factors of $2\pi$ and $4\pi$. The structure factor of \cref{eq:4.2-Sab} is per volume, in \si{eV^2} in natural units; a table per cell holds $\Omega S$, in \si{eV^{-1}}, with $\Omega = \SI{40.05}{\angstrom^3}$ for silicon; a rate per target mass divides by $\rho_T$ in the halo prefactor. Spin is summed. A table of loss functions is turned into structure factors by \cref{eq:4.4-SL}, $S = (2q^2/e^2)\,L$ at zero temperature with $e^2 = 4\pi\alpha$; in Gaussian units the same symbol $e^2$ means $\alpha$ and the factor changes by $4\pi$. A dielectric table must say which function it lists: the head $\varepsilon_{00}$ of the dielectric matrix, whose loss function $\varepsilon_2/|\varepsilon_{00}|^2$ is the random-phase value without local fields, or the macroscopic function $\varepsilon_{\mathrm{mac}} = 1/(\varepsilon^{-1})_{00}$, whose loss function includes them. The two loss functions differ by the local-field step, 9 percent at the node of \cref{we:21-libentry}. Momenta are unreduced (\cref{app:A.3}); a table that lists $|\vect{q} + \vect{G}|$ for the points of its calculation mixes directions at equal length, and an average over bins of $|\vect{q}|$ is then an average over directions, which the table states.

The row basis and its index order are declared with the matrix. The contraction is $\sum_{ab}\mathfrak{D}_{ab}S_{ab}$ with the rows of \cref{tab:10-minimal}, and the order of the indices of $R^0_{\vect{G}\vect{G}'}$ fixes which factor is conjugated. Transposing $R^0_{\vect{G}\vect{G}'}$ conjugates the cross terms, which is immaterial in silicon, where every element is real, but not in GaAs (\cref{sec:12.3}). The final-state tables add the labeling of pairs (retarded) and the convention of the sector split, and a table of Migdal responses states its energy axis, the total electronic energy transfer $\omega_e$ (\cref{sec:12.1}).

The truncation label is written out in all eight components of \cref{eq:8.5-label}, not in the deviation notation that suits running text, together with the reference state (normal, superconducting, spin-wave), the screening set and the split convention. The label leaves out the exchange--correlation functional, the pseudopotentials, the scissor and the grids; these approximations are measured by tests and convergence (\cref{chap:23,chap:24,chap:27,chap:25}), and they go into the material inputs with the results of those tests. The scissor gap and the pair-conversion gap need separate entries. The gap to which a scissor shifts the bands and the gap with which a cascade or a step model converts energy into pairs are different numbers in general, and a rate that uses both inherits two thresholds. Near \SI{1}{eV} the energy nodes of the tables of \cref{we:21-libentry} are \SI{0.22}{eV} apart, more than the \SI{0.06}{eV} between the room-temperature and low-temperature gaps of silicon (\cref{sec:10.10}), so the gap of such a table cannot be read off its nodes and has to be stated. The tables of \cref{we:21-libentry} state a scissor to the measured gap at zero temperature, \SI{1.17}{eV}; the \SI{1.11}{eV} that accompanies them in their package is the gap of the step model of pair creation (\cref{sec:10.10}). The interpolation rule between nodes is part of the table too: near \SI{10}{eV} the same loss function changes by 10 percent from one energy node to the next.

\subsection{Versions and provenance}\label{sec:21.6-versions}

Every array of a library is identified by a hash of its content. It carries the list of its parents: the structure, the pseudopotentials, the code and its version, the input files, and for a derived table the tables it was computed from. Two arrays with the same name and different content can then not be confused, and a limit can name the exact arrays it used. Versions follow one rule for each kind of change. A change of a truncation label or of a convention changes rates and makes a new major version. Added tables, such as a new row pair, a new material or a new uncertainty template, change no existing rate and make a minor version. A numerical refinement, a finer grid or a larger cutoff, is a patch only if no tabulated value moves by more than the convergence tolerance the table declares. A larger move means the old values were outside their stated accuracy, and it makes a major version as well. A published version is never changed and carries a persistent identifier, so a limit cites the library the way it cites a paper.

A version is released with its checks, each stored as a value rather than a verdict. The $f$-sum rule \cref{eq:4.3-fsum} holds to \SI{2}{\percent} below the cutoff of the bands, and a spectrum that ends early stores the fraction of the $f$-sum it holds at each $|\vect{q}|$. Positivity, \cref{eq:4.2-positivity}, is checked through the smallest eigenvalue of every stored matrix $S_{ab}$, with a tolerance stated by the library, for instance $-10^{-6}$ of its trace; the Kramers--Kronig residual is stored likewise, for instance below \SI{1}{\percent}, and the two signs of $\omega$ satisfy detailed balance. In silicon and germanium the symmetry-forbidden parts of the density--current and current blocks stay below $10^{-3}$ of the allowed parts. The cascade tables sum to one and are nonnegative, their derivatives reproduce the moment recursion to $10^{-9}$, and below the first impact ionization, at $(1 + \theta_{\mathrm{ii}})\Eg = \SI{2.32}{eV}$ in silicon (\cref{sec:18.3}), they give exactly one pair. The step columns never decrease. Where optical, energy-loss or inelastic x-ray measurements exist, the comparison with them is stored in the same way. A table that misses a check by a few percent may still serve an analysis that needs less, and only the stored value lets its user decide.

A limit published with a library names the library and its version, the truncation label of each table it used and the cascade model, and gives per bin the count, the exposure and the detector response. The limit itself is then reproduced from the tables. Its theory band follows from the uncertainty tables and their keys (\cref{sec:21.4}). A recast for another coupling, another halo or a corrected response evaluates the kernel \cref{eq:21.1-kernel} again instead of rescaling a published curve, which \cref{sec:21.5} shows to be exact only under a condition on the bins. The guide to targets of \cref{sec:21.7} reads its inputs from the same records: the loss function $L(\vect{k},\omega)$ of each candidate and its elastic constants. The provenance it reports for each figure of merit, measured, computed or estimated, is the provenance stored with those tables.

\begin{workedexample}{The silicon loss function at one table node}\label{we:21-libentry}
Knapen, Kozaczuk and Lin computed the random-phase dielectric function of silicon from Kohn--Sham states~\cite[Sect.~III\,B]{knapen2021dark}. They used the TB09 functional, the meta-generalized-gradient functional of Tran and Blaha (2009) chosen for band gaps, a plane-wave cutoff of \SI{500}{eV}, an $8\times8\times8$ grid of $\vect{k}$ points, 70 bands per spin and a scissor to the measured gap. The tables distributed with their package~\cite{knapen2022darkelf} list $\varepsilon_1$ and $\varepsilon_2$ against $\omega$ and $|\vect{q} + \vect{G}|$, in one file for the head $\varepsilon_{00}$ and in one for $\varepsilon_{\mathrm{mac}}$. We write the entry at the node nearest to $(q_0,\SI{10}{eV})$, with $q_0 = \alpha m_e = \SI{3728.9}{eV}$.

\emph{Node.} $|\vect{q}| = \SI{3732.0}{eV} = 1.0008\,q_0$ and $\omega = \SI{9.8585}{eV}$. There the head is $\varepsilon_{00} = 1.838 + 0.591\,\ii$ and the macroscopic function $\varepsilon_{\mathrm{mac}} = 1.763 + 0.482\,\ii$.

\emph{Value.} With local fields, $L = \Imag(-1/\varepsilon_{\mathrm{mac}}) = 0.1442$. By \cref{eq:4.4-SL}, $2q^2/e^2 = \SI{3.038e8}{eV^2}$, so $S = \SI{4.380e7}{eV^2}$ per volume and $\Omega S = \SI{0.2283}{eV^{-1}}$ per cell. The same node gives $L = \varepsilon_2/|\varepsilon_{00}|^2 = 0.1585$ without local fields and, from the head's $\varepsilon_2 = 0.591$, the independent-particle $S^{\mathrm{IP}} = \SI{1.796e8}{eV^2}$.

\emph{Label.} $\mathfrak{T} = (n_B = 0,\, n_g = 0,\, n_A = 0,\, \mathcal{X} = \mathrm{RPA + LF},\, n_T = 0;\ \mathcal{K},\, \mathfrak{M},\, n_D$ not applicable$)$, with $T = 0$, a normal reference state, external and local-field screening, and no sector split, since a table at $n_g = 0$ has no phonon-dressed part. Off the label: TB09, the cutoff, the grid, the bands and the scissor, as above. The package keeps a separate file of material constants, with $\Eg = \SI{1.11}{eV}$ and $\varepsilon_{\mathrm{eh}} = \SI{3.6}{eV}$ for the conversion of energy into pairs by the step model, and an entry of the library would carry that gap beside the gap of the scissor.

\emph{Uncertainty.} The local-field step is computed, so it moves the central value: $\delta_{\mathcal{X}} = 0.1442/0.1585 - 1 = -0.090$ at this node, a tilt by \cref{eq:8.8-fsum}. The theory uncertainty that remains is the excitonic correction. No calculation at finite momentum gives it for silicon, so the node carries the flag of a missing template. \Cref{sec:10.10} estimates the step at $+0.28$ over the energies a \SI{3}{MeV} particle reaches, from a model at \SI{0.91}{keV}, and a band takes the estimate from \cref{sec:10.10} with that basis named.

\emph{Checks.} At this $|\vect{q}|$ the table holds 0.916 of the valence $f$-sum $(\pi/2)\omega_p^2$, $\omega_p = \SI{16.6}{eV}$, below its last energy, \SI{89.2}{eV}; the missing weight lies above about \SI{70}{eV}, which the 70 bands do not reach~\cite[Fig.~1]{knapen2021dark}. The three other nodes within \SI{50}{eV} of $|\vect{q}|$, at 3689, 3722 and \SI{3776}{eV}, give $L = 0.1431$, 0.1402 and 0.1445. The four nodes span 3.0 percent, partly from direction and partly from $|\vect{q}|$, and the entry states that its direction average is a bin average over such nodes. The next energy node, \SI{10.26}{eV}, gives $L = 0.1295$, so at \SI{10}{eV} linear interpolation gives $L = 0.1390$ and $S = \SI{4.22e7}{eV^2}$.

\emph{Comparison.} The empirical-pseudopotential model of \cref{we:8-pairs} gives, at $q_0$ along [100] and \SI{10}{eV}, $S^{\mathrm{IP}} = \SI{1.7e8}{eV^2}$ against \SI{1.74e8}{eV^2} here, and $S = \SI{4.89e7}{eV^2}$ with local fields against \SI{4.22e7}{eV^2}, a ratio of 0.864. The screening factor $S/S^{\mathrm{IP}}$ is 0.29 in the model and 0.242 in the tables, and the local-field step $-0.06$ against $-0.09$. The two calculations carry the same label. They differ in what is off the label: a model potential with a gap of \SI{0.82}{eV} against TB09 with a scissor, and one direction against a mixture. Only the provenance of the table tells a user which of the two a limit used.
\end{workedexample}
